% Drawing 138/06-20 -- folder 138, batch 06, archive page 113 (PDF page 114).
% Transcription: pour a) : une croix ; une boucle traversée d'un segment ; deux boucles accolées ; une boucle avec un segment intérieur (ces deux dernières séparées par un trait courbe). Pour b) : un segment suivi d'une fourche ; des esquisses biffées ; « rigide ».
% Kind: small maps (degree 4). Method 2 (schematic redraw from the facsimile). Status: ok.
% Notes: For a), above the long curved stroke: a cross (two edges crossing at a vertex); a loop at a vertex crossed by a segment (a stroke in from the left and one up, at the same vertex); two loops joined at a vertex; a loop with an inner segment (the segment enters at the top vertex and ends inside). The long curved stroke that separates a) from b) runs under the last three. For b): a segment followed by a fork; « rigide ». Every edge is drawn between the same vertex points. Struck and omitted: the heavily scribbled sketches of b) (two or three small maps under zigzag strokes, between the fork and « rigide ») and the scribble under « rigide »; the earlier redraw drew them as grey outlines. Vertex positions and loop shapes read on the crop (to about one unit). Small planar maps listed by degree (number of edges) in the text; redrawn, not encoded (each would be a map/1 with at most a few edges; the text gives the classification). Redrawn 2026-11-09 (issue 128): to check again.
% Coordinates: 1 unit = 1 mm = 1/100 of the crop width (bbox in page fractions: [0.52, 0.685, 0.9, 0.87]).
% Redrawn by Claude (claude-opus-5-5), issue #73, 2026-11-09; not checked by a human; not compiled (no TeX in the session).
% Rights: reproduces part of the work -- behind RENDER_TRANSCRIPTS, never published by itself.
\documentclass[tikz,border=3pt]{standalone}
\usepackage{amsmath,amssymb}
\usetikzlibrary{patterns,arrows.meta}
\tikzset{pencil/.style={black!45}, struck/.style={black!35}, mark/.style={circle,draw,inner sep=1.3pt}, pt/.style={circle,fill,inner sep=0.9pt}, opt/.style={circle,draw,fill=white,inner sep=1.1pt}}
\begin{document}
\begin{tikzpicture}[x=1mm,y=1mm,line cap=round,line join=round,font=\scriptsize]
\draw (36.3,-7.1) -- (43.3,-11.5) -- (48.5,-16.2);
\draw (53.5,-9.3) -- (43.3,-11.5) -- (35,-16.5);
\node[pt] at (43.3,-11.5) {};
\draw (66.1,-12) -- (75.1,-12.2);
\draw (75.1,-12.2) -- (76.3,-3.7);
\draw plot[smooth, tension=0.55] coordinates {(75.1,-12.2) (71.9,-14.7) (71.7,-17.9) (72.9,-21.8) (80,-23.6) (85.5,-20.2) (86.4,-13.9) (82.4,-10.6) (77.6,-11.2) (75.1,-12.2)};
\node[pt] at (75.1,-12.2) {};
\draw plot[smooth, tension=0.55] coordinates {(54.1,-27.3) (55.5,-23.5) (51.5,-21.8) (45.2,-23) (43.3,-27.3) (46,-30.1) (51.5,-29.3) (54.1,-27.3)};
\draw plot[smooth, tension=0.55] coordinates {(54.1,-27.3) (57.1,-25) (62.6,-24.2) (65.8,-27.3) (65,-32.8) (57.8,-35.3) (54,-32.8) (54.1,-27.3)};
\node[pt] at (54.1,-27.3) {};
\draw plot[smooth, tension=0.55] coordinates {(87.5,-27.3) (85.4,-28.6) (84.8,-32) (86.5,-35.7) (89.8,-36.8) (93.1,-34.6) (94.2,-31.3) (92.6,-28) (89.8,-26.9) (87.5,-27.3)};
\draw (89.9,-20.2) -- (87.5,-27.3) -- (91,-32);
\node[pt] at (87.5,-27.3) {};
\draw plot[smooth, tension=0.55] coordinates {(6.6,-21.8) (18.4,-21) (27.9,-23.5) (34.2,-28.1) (42.2,-32.8) (57.8,-36.8) (73.7,-38.4) (90.7,-38.4)};
\draw (1.8,-41) -- (14.5,-41);
\draw (14.5,-41) -- (23.1,-36.8);
\draw (14.5,-41) -- (21.4,-45.3);
\node[pt] at (1.8,-41) {};
\node[pt] at (14.5,-41) {};
\node[inner sep=1pt] at (70.4,-60.4) {\small rigide};
\end{tikzpicture}
\end{document}
